\documentclass[10pt,a4paper]{amsart}
\usepackage[margin=19mm]{geometry}
\usepackage{amsmath,amssymb,mathtools}
\usepackage[scr=rsfs,cal=euler]{mathalfa}
\usepackage{xfrac}
\usepackage[unicode,hidelinks,bookmarks=false]{hyperref}
\newcommand{\ie}{i.e.\ }
\newcommand{\cf}{cf.\ }
\newcommand{\resp}{respectively\ }
\newcommand{\Subsection}{\subsection}
\providecommand{\nobreakdash}{\nobreak\hbox{-}}
\setlength{\emergencystretch}{2em}
\multlinegap0pt
\usepackage{amssymb}
\usepackage{mathtools}
\usepackage[shortlabels]{enumitem}
\newtheorem{lemma}{Lemma}
\newtheorem{corollary}{Corollary}
\newtheorem{proposition}{Proposition}
\usepackage{cleveref}
\crefname{lemma}{Lemma}{Lemmas}
\crefname{corollary}{Corollary}{Corollarys}
\crefname{proposition}{Proposition}{Propositions}
\newcommand{\GG}{\mathbb{G}}
\newcommand{\PP}{\mathbb{P}}
\DeclareMathOperator{\Sing}{Sing}
\newcommand{\ZZ}{\mathbb{Z}}
\DeclareMathOperator{\id}{id}
\title{Kodaira-type classification of singular fibers of some minimal abelian fibrations}
\author{Yoon-Joo Kim}
\date{Source: arXiv:2603.02347v1 (CC BY 4.0)}
\begin{document}
\maketitle

\noindent\emph{Source and licence.} Kodaira-type classification of singular fibers of some minimal abelian fibrations. Version arXiv:2603.02347v1 (CC BY 4.0), distributed by arXiv under the Creative Commons Attribution 4.0 International licence, http://creativecommons.org/licenses/by/4.0/, as recorded in the arXiv licence metadata for this exact version and shown on the abstract page https://arxiv.org/abs/2603.02347v1. Full licence notice: \texttt{licenses/arxiv-2603.02347-CC-BY-4.0.txt}.\par\smallskip

\noindent\textit{Editorial scope.} Counted unit: the article's lemma iii:psi-action (source0.tex lines 2360--2363). The article's complete original proof of it is delivered with it (source0.tex lines 2364--2379, one \texttt{proof} environment of 16 lines). No mathematics is written, repaired or re-derived: the statement and the proof are the article's own lines, in the article's own notation, and no retained line is substituted or re-flowed.  Retained uncounted as the source-local context the proof needs: lines 1322--1324; lines 1564--1566; lines 1667--1673; lines 2353--2358; lines 2950--2956.  Nothing was omitted from any retained span, so the package carries no omission record.  No retained line contains a citation, so the wrapper carries no bibliography and no bibliography entry.  No retained line contains a figure environment, an image-inclusion command or drawing code, so no image is delivered and the package declares no asset dependency.  Editorial formatting only, in addition: an AMS \texttt{amsart} preamble that restates the article's own declarations and macros used by the retained lines, namely the environments \texttt{lemma}, \texttt{corollary}, \texttt{proposition} on separate counters, together with the standard packages those lines need.  The retained environments print in this document as Lemma 1, Corollary 1, Proposition 1 in order of appearance, whereas the article prints the same items with the numbering of its own full document.  The complete native source of the article as distributed in the arXiv e-print for this exact version is bundled unchanged as \texttt{sources/195611/source0.tex}.
		\begin{lemma} \label{lem:stabilizer of Ir}
			Assume $f : X \to S$ has type $\mathrm I_r$ for $r \ge 1$. Then every nontrivial stabilizer of the $P_s$-action on $X_s$ is $\GG_m$.
		\end{lemma}

			\begin{corollary} \label{cor:pi_0-action on X_s}
				Assume $f$ is strictly semistable of type $\mathrm I_r$. Then $\pi_0(P_s) \cong \ZZ/r$ acts on $X_s = \sum_{i=1}^r E_i$ by cycling the components $E_i$.
			\end{corollary}

			\begin{proposition} \label{prop:varphi-action}
				\begin{enumerate}
					\item There exists a natural automorphism $\varphi$ of order $d$ on $Y$ making $g : Y \to T$ equivariant.
					\item The $Q_t$-and $\langle \varphi \rangle$-actions on $Y_t$ are related by
					\[ \varphi(a \cdot y) = \psi(a) \cdot \varphi(y) \qquad\mbox{for}\quad a \in Q_t , \ y \in Y_t .\]
				\end{enumerate}
			\end{proposition}

			\begin{corollary} \label{cor:classification of multiple fiber iii:varphi-action on Sing}
				\begin{enumerate}
					\item The $\langle \varphi \rangle$-action on $\Sing(Y_t) = \bigsqcup_{i=1}^{kr} (F_i \cap F_{i+1})$ is a translation by $a \in A_s$ on every component.
					\item $\tilde E_i \to A_s/a$ is a free $\varphi$-quotient of $\tilde F_i \to A_s$. In particular, it is a $\PP^1$-bundle.\qed
				\end{enumerate}
			\end{corollary}

		\begin{lemma} \label{lem:pi_0 sequence splits}
			Let $G$ be a commutative locally algebraic group.
			\begin{enumerate}
				\item If $\pi_0(G)$ is finitely generated, then $0 \to G^{\circ} \to G \to \pi_0(G) \to 0$ splits.
				\item If $\pi_0(G)$ is finite and $G^{\circ}$ is unipotent, then the splitting is unique.
			\end{enumerate}
		\end{lemma}

			\begin{lemma} \label{lem:classification of multiple fiber iii:psi-action}
				There exists a primitive $k$-th root of unity $\zeta \in \GG_m \subset Q_t^\circ$ such that the inertia $\langle \psi \rangle$-action on $Q_t$ is
				\[ \psi(f_i) = \zeta^i \cdot f_i \qquad\mbox{for}\quad f_i \in Q_t \mbox{ with } \pi_0(f_i) = i \ \in \ \pi_0(Q_t) = \ZZ/kr .\]
			\end{lemma}

			\begin{proof}
				We may write $Q_t = \ZZ/kr \times Q_t^\circ$ by \Cref{lem:pi_0 sequence splits}, and write its element by $(i, g) \in \ZZ/kr \times Q_t^\circ$. Recall the equality $Q_t^\psi = P_s$ and $\pi_0(P_s) = k \cdot \ZZ/kr \subset \ZZ/kr = \pi_0(Q_t)$, so we already have $\psi(i,g) = (i,g)$ when $k \mid i$. Recall the equality $\varphi ((i,g) \cdot y) = \psi(i,g) \cdot \varphi(y)$ of \Cref{prop:varphi-action}. Since $\varphi$ preserves all $F_i$ by assumption and $\pi_0(Q_t) = \ZZ/kr$ cycles the components $F_i$ (\Cref{cor:pi_0-action on X_s}), this shows that $\psi$ is of the form
				\[ \psi(i,g) = (i, \psi_i (g)) \qquad\mbox{for}\quad \psi_i : Q_t^\circ \to Q_t^\circ .\]
				Since $\psi$ acted trivially on $Q_t^\circ$, we have $\psi(i,g) = \psi(0,g) \cdot \psi(i,1) = (0,g) \cdot (i,\psi_i(1)) = (i, \psi_i(1) \cdot g)$. This shows $\psi_i(g) = \zeta_i \cdot g$ for $\zeta_i = \psi_i(1) \in Q_t^\circ$. The equality $\psi(i,g^i) = \psi(1, g)^i$ shows $\zeta_i = \zeta^i$ for a single $\zeta \in Q_t^\circ$. Furthermore, $\zeta$ is $k$-torsion because $\psi^k = \id$. Again $Q_t^\psi$ is precisely $P_s$, so this implies that $\zeta$ has order precisely $k$.
				
				\medskip
				
				It remains to prove that $\zeta \in Q_t^\circ$ is contained in $\GG_m$. Take a point $y \in F_1 \cap F_{kr}$, a unique rational curve $R \subset F_1$ through $y$, and a unique point $y' = R \cap F_2 \subset F_1 \cap F_2$. Recall from \Cref{lem:stabilizer of Ir} that $Q_t$ acts on $\Sing(Y_t)$ transitively by cycling the components (with stabilizer $\GG_m$). Hence there exists $f_1 \in Q_t$ such that $\pi_0(f_1) = 1$ and $y' = f_1 \cdot y$.
				
				Now apply $\varphi$ to $y, y' \in R \subset F_1$ to yield two points on a rational curve $\varphi(y), \varphi(y') \in \varphi(R) \subset F_1$. By \Cref{cor:classification of multiple fiber iii:varphi-action on Sing}, if we fix any $f_0 \in Q_t^\circ$ whose image under $Q_t^\circ \to A_s$ is $a$, then $\varphi(y) = f_0 \cdot y$ and $\varphi(y') = f_0 \cdot y'$. We now have a sequence of identities
				\begin{align*}
					\varphi(y') &= \varphi(f_1 \cdot y) = \psi(f_1) \cdot \varphi(y) = (\zeta \cdot f_1) \cdot (f_0 \cdot y) \\
					&= \zeta \cdot (f_0 \cdot (f_1 \cdot y)) = \zeta \cdot (f_0 \cdot y') = \zeta \cdot \varphi(y') .
				\end{align*}
				But $\varphi(y') \in F_1 \cap F_2$ has a stabilizer $\GG_m$ again by \Cref{lem:stabilizer of Ir}. This means $\zeta \in \GG_m \subset Q_t^\circ$.
			\end{proof}

\end{document}
